\section{Imagens Docker}

\begin{frame}
  \frametitle{O que é uma Container Image?}
  Pacote padronizado com tudo o necessário para rodar um container:

  \vspace{0.5em}
  \begin{itemize}
    \item \textbf{PostgreSQL}: Binários do banco, configurações e dependências
    \item \textbf{Python}: runtime, código da aplicação e suas dependências
  \end{itemize}
\end{frame}

\begin{frame}
  \frametitle{Características das Imagens}
  \begin{itemize}
    \item \textbf{Imutabilidade}: Uma vez criada, não pode ser modificada;
          apenas cria-se uma nova imagem ou adiciona-se camadas sobre uma existente
    \item \textbf{Camadas (\textit{Layers})}: Cada camada representa um conjunto
          de mudanças no sistema de arquivos: adição, remoção ou modificação de arquivos
  \end{itemize}

  \vspace{0.5em}
  Essas características permitem \textbf{estender imagens existentes}, por exemplo,
  partir de uma imagem base Python e adicionar camadas com dependências e código.
\end{frame}

\begin{frame}
  \frametitle{Encontrando Imagens — Docker Hub}
  Marketplace padrão para armazenar e distribuir imagens. Categorias:

  \vspace{0.5em}
  \begin{itemize}
    \item \textbf{Docker Official Images}: Repositórios curados pelo Docker
    \item \textbf{Docker Hardened Images}: Imagens mínimas e seguras, praticamente zero CVEs
    \item \textbf{Docker Verified Publishers}: Imagens de fornecedores verificados
    \item \textbf{Docker-Sponsored Open Source}: Mantidas por projetos open source patrocinados
  \end{itemize}
\end{frame}

\begin{frame}
  \frametitle{Comandos Básicos para Imagens}
  Buscar uma imagem:

  \vspace{0.3em}
  \texttt{docker search docker/welcome-to-docker}

  \vspace{0.5em}
  Baixar a imagem (cada linha = uma camada):

  \vspace{0.3em}
  \texttt{docker pull docker/welcome-to-docker}

  \vspace{0.5em}
  Listar imagens locais:

  \vspace{0.3em}
  \texttt{docker image ls}

  \vspace{0.5em}
  Ver as camadas de uma imagem:

  \vspace{0.3em}
  \texttt{docker image history docker/welcome-to-docker}
\end{frame}